\subsection{观测数量与污染水平的预登记敏感性分析}
\label{sec:sensitivity-protocol}

为补充单一压力点的评价，另以独立脚本\texttt{sensitivity.py}研究适用范围。
在生成数据前保存的\texttt{research/sensitivity\_plan.json}固定根种子980221，
观测数$n\in\{8,12,16,24\}$、名义污染比例
$\epsilon\in\{0,0.1,0.25,0.4\}$，以及静态异常、恒速异常两个场景。
每个网格单元独立生成12例，合计32格、384例。各格使用由场景编号、观测数和
污染水平派生的独立伪随机子流；三种方法在每格使用完全相同的案例，因此对照差异
按案例配对。固定比较方法仅为最终冻结方案、恒速稳健基线和静态加权最小二乘，
总计1152次预测，不依此重新选择方法或修改参数。

脚本先验证研究冻结、关键源码与开发证据，再生成数据。协议保存预登记文件、
本脚本及冻结方法的摘要；这些本地记录用于审计顺序，不等于第三方时间戳认证。
每格报告平均位置误差、经验CEP90、计算工作量、壁钟时间和实际模型选择频率，
同时保存逐例输入、预测与CV诊断。污染数采用生成器的
$\operatorname{round}(n\epsilon)$，因而小样本下的实际比例具有离散性；
结果必须并列记录名义比例与实际比例。

观测数同时影响信息量与CV门控。默认三折时间交错划分在$n=8$时含训练样本数
为5的折，低于运动模型要求的6份训练观测；若冻结方案采用该CV配置，运动模型会
因可用性条件被排除。该机制可由完成折数与运动可用标记核对，不能将其误读为数据
支持目标静止。若最终方案不含CV，则相关字段记为不适用，不补造门控结果。
本分析仅描述预设合成网格，每格12例不足以精确定位稳健性相变；不作逐格显著性
或多重比较确证，不重新校准误差区域，也不提出新的覆盖率保证。
